% design.tex -- the 13 slides of design/thbtalk-preview.html with the same content, for test/design_check.py
\documentclass[aspectratio=169,10pt]{beamer}
\usepackage[utf8]{inputenc}
\usepackage[english]{babel}
\usepackage{thbtalk}
\title{Adaptive Sampling for Sparse Sensor Networks}
\subtitle{A field study on the Upper Rhine}
\author{A.~Musterfrau}
\institute{Technische Hochschule Bingen}
\date{19 September 2026}
\thbvenue{ICSE Workshop, Mainz}
\thblogotext{TH Bingen}
\thbcontact{a.musterfrau@th-bingen.de \textperiodcentered\ doi.org/10.0000/xyz}
\thboutlinenote{1}{Where static sampling breaks down in the field}
\thboutlinenote{2}{A variance-driven scheduler with a submodular guarantee}
\thboutlinenote{3}{Eleven weeks on the Upper Rhine, 18 nodes}
\begin{document}
\thbtitleframe
\thboutline
\section{Motivation}
\begin{frame}{Why sparse networks fail}{Three observations from two deployments}
  \begin{itemize}
    \item Node dropout is bursty, not uniform.
    \item Static sampling wastes 40\,\% of the energy budget.
    \item Reconstruction error grows superlinearly below 12 active nodes.
  \end{itemize}
  \begin{takeaway}
    Sampling rate should follow measured field variance, not the clock.
  \end{takeaway}
\end{frame}
\begin{frame}{The gap in the literature}
  \thbquote{Placement is treated as a one-off design decision, although the field
    it samples is non-stationary on the scale of days.}{Krause \& Guestrin, JMLR 2008}
\end{frame}
\section{Method}
\begin{frame}{Problem statement}
  We minimise the expected reconstruction error under an energy constraint:
  \begin{keyeq}
    \begin{equation}
      \min_{\pi}\; \mathbb{E}\!\left[\lVert x - \hat{x}_{\pi}\rVert_2^2\right]
      \quad\text{s.t.}\quad \sum_{t} c(\pi_t) \le E_{\max}.
    \end{equation}
  \end{keyeq}
  \begin{theorem}[Convex relaxation]
    For submodular $c(\cdot)$ the relaxed problem admits a $(1-1/e)$ approximation.
  \end{theorem}
\end{frame}
\begin{frame}{Algorithm}
  \begin{pseudocode}{Variance-driven scheduler}
    \State $S \gets \emptyset$
    \While{$\text{budget}(S) < E_{\max}$}
      \State $v^\ast \gets \arg\max_{v \notin S} \Delta(v \mid S)$
      \State $S \gets S \cup \{v^\ast\}$
    \EndWhile
    \State \Return $S$
  \end{pseudocode}
\end{frame}
\begin{frame}[fragile]{Reference implementation}
  \begin{code}[language=Python]
def schedule(nodes, budget):
    active = []
    while cost(active) < budget:
        best = max(nodes, key=lambda v: gain(v, active))
        active.append(best)   # greedy, submodular gain
    return active
  \end{code}
\end{frame}
\section{Results}
\begin{frame}{Field data}
  \begin{twocols}
    \thbplaceholder[1.0]{89.76pt}{sensor map}
    \colbreak
    \begin{itemize}
      \item 18 nodes, 11 weeks
      \item 2.1\,M samples
      \item 3 dropout events
    \end{itemize}
    Adaptive scheduling kept coverage above 92\,\% throughout.
  \end{twocols}
\end{frame}
\thbfullfigure{stripes.png}{Deployment area, 11 weeks of continuous operation. Nodes in green were rescheduled at least once.}
\begin{frame}{Headline result}
  \vfill
  \thbnumber{39\,\%}{lower reconstruction error than fixed-rate sampling at the same energy budget}
  \vfill
\end{frame}
\begin{frame}{Comparison}
  \begin{center}
    \begin{tabular}{lrrr}
      \toprule
      \thbhead{Method} & \thbhead{RMSE} & \thbhead{Energy} & \thbhead{Uptime} \\
      \midrule
      Fixed rate       & 0.72 & 3.0 & 88\,\% \\
      Round robin      & 0.58 & 3.0 & 90\,\% \\
      Adaptive (ours)  & \textbf{0.44} & 3.0 & \textbf{92\,\%} \\
      \bottomrule
    \end{tabular}
  \end{center}
  {\centering\footnotesize\color{thbmuted} Averaged over 11 weeks; lower is better.\par}
\end{frame}
\begin{closing}
  \item Sampling should follow field variance, not the clock.
  \item 39\,\% lower error at an unchanged energy budget.
  \item Code and the 11-week dataset are public.
\end{closing}
\end{document}
